\hypertarget{exception-groups-and-traceback-enhancements}{%
\chapter{EXCEPTION GROUPS AND TRACEBACK
ENHANCEMENTS}\label{exception-groups-and-traceback-enhancements}}

\hypertarget{exception-groups-exceptiongroup-baseexceptiongroup-pep-654}{%
\subsection{\texorpdfstring{19.1 Exception Groups
(\texttt{ExceptionGroup} \& \texttt{BaseExceptionGroup} / PEP
654)}{19.1 Exception Groups (ExceptionGroup \& BaseExceptionGroup / PEP 654)}}\label{exception-groups-exceptiongroup-baseexceptiongroup-pep-654}}

Introduced in Python 3.11, \passthrough{\lstinline!ExceptionGroup!} and
\passthrough{\lstinline!BaseExceptionGroup!} enable raising and handling
multiple unrelated exceptions simultaneously. This is critical for
concurrent frameworks (such as \passthrough{\lstinline!asyncio!}
TaskGroups) where multiple background operations can crash concurrently.

\hypertarget{exception-group-class-layout}{%
\subsubsection{1. Exception Group Class
Layout}\label{exception-group-class-layout}}

An \passthrough{\lstinline!ExceptionGroup!} wraps a descriptive string
and a list of sub-exceptions:

\begin{lstlisting}[language=Python]
class ExceptionGroup(BaseException):
    def __init__(self, message: str, exceptions: Sequence[BaseException]) -> None:
        self.message = message
        self.exceptions = list(exceptions)
\end{lstlisting}

\begin{center}\rule{0.5\linewidth}{0.5pt}\end{center}

\hypertarget{the-except-clause-and-exception-tree-filtering}{%
\subsection{\texorpdfstring{19.2 The \texttt{except*} Clause and
Exception Tree
Filtering}{19.2 The except* Clause and Exception Tree Filtering}}\label{the-except-clause-and-exception-tree-filtering}}

To handle individual branches of an exception group, Python 3.11
introduced the \passthrough{\lstinline!except*!} statement. Unlike a
traditional \passthrough{\lstinline!except!} statement, which catches a
single exception instance, \passthrough{\lstinline!except*!} extracts a
matching subset from the Exception Group hierarchy, letting unmatched
exceptions propagate.

\hypertarget{execution-trace-of-exception-splitting}{%
\subsubsection{1. Execution Trace of Exception
Splitting}\label{execution-trace-of-exception-splitting}}

Consider the following exception structure and matching logic:

\begin{lstlisting}
            ExceptionGroup("Main Group")
            /                        \
    ValueError("Error A")       TypeError("Error B")
\end{lstlisting}

\begin{lstlisting}[language=Python]
try:
    raise ExceptionGroup("Main Group", [ValueError("Error A"), TypeError("Error B")])
except* ValueError as eg:
    print("Caught Val:", eg.exceptions)
\end{lstlisting}

During execution, the runtime: 1. Catches the
\passthrough{\lstinline!ExceptionGroup!}. 2. Filters out the
\passthrough{\lstinline!ValueError("Error A")!} exception and wraps it
in a new sub-ExceptionGroup. 3. Binds this sub-ExceptionGroup to the
local variable \passthrough{\lstinline!eg!}. 4. Leaves the
\passthrough{\lstinline!TypeError("Error B")!} exception in the original
group, which continues propagating up the exception handler chain.

\begin{center}\rule{0.5\linewidth}{0.5pt}\end{center}

\hypertarget{traceback-representation-and-add_note}{%
\subsection{\texorpdfstring{19.3 Traceback Representation and
\texttt{add\_note()}}{19.3 Traceback Representation and add\_note()}}\label{traceback-representation-and-add_note}}

\begin{itemize}
\tightlist
\item
  \textbf{Traceback Trees}: Because exception groups represent
  hierarchical trees of errors, CPython's traceback generator is
  modified to format tracebacks as tree diagrams.
\item
  \textbf{PEP 678 Exception Notes}: The
  \passthrough{\lstinline!add\_note(note)!} method allows adding custom
  text strings directly to exceptions. These notes are appended to the
  exception's traceback output when printed, aiding debugging without
  modifying the exception's arguments.
\end{itemize}

\begin{center}\rule{0.5\linewidth}{0.5pt}\end{center}
